\begin{abstractzh}
学位论文的规范排版是准确表达研究内容和组织学术材料的重要基础。本示例围绕贵州大学硕士学位论文的编排要求，展示从题名页、答辩委员会名单和声明页，到中英文摘要、目次、正文及结尾部分的完整组织方式。文档中的姓名、学号、单位及研究内容均为示例，使用时需要由作者据实替换。

模板将论文基本信息集中设置，统一生成前置页面，并分别控制正文、章节标题、图表题注和参考文献的字体与字号。针对中文论文中常见的长网址、多行表格、跨页公式和附录编号问题，示例给出了可直接使用的写法。正文后的致谢、参考文献、附录、索引、作者简历和后记采用独立的样式，同时保持页码连续。

本摘要仅说明模板的用途，不报告任何实际科研发现。正式撰写时，中文摘要应概括研究目的、方法、结果和结论，突出研究工作的创新或改进，通常控制在三百至一千字，并选择三至八个关键词。英文摘要应与中文摘要的内容对应。作者应在提交前逐项核对论文信息、声明签名、参考文献及院系补充要求。
\end{abstractzh}

\begin{abstracten}
Consistent thesis formatting supports the clear presentation of research and the organization of academic material. This example demonstrates the structure of a master's thesis at Guizhou University, from the title page, defense committee and declarations to the Chinese and English abstracts, contents, main text and end matter. Names, student details, affiliations and research descriptions are illustrative and must be replaced with accurate information.

The template stores thesis information centrally and generates the preliminary pages with consistent fonts and sizes. Examples address long URLs, multiline table cells, equations spanning pages and appendix numbering. Acknowledgements, references, appendices, an index, the author's curriculum vitae and a postscript have separate styles while retaining continuous page numbers.

This abstract describes the template and reports no research findings. In an actual thesis, the abstract should state the research purpose, methods, results and conclusions, with emphasis on the contributions or improvements. The Chinese abstract normally contains 300 to 1,000 Chinese characters and three to eight keywords. The English abstract should convey the same content. Authors should verify all personal information, declaration signatures, references and additional departmental requirements before submission.
\end{abstracten}
